\documentclass{article}
\usepackage{amsmath,amssymb}
\usepackage{csquotes}
\usepackage{epigraph}
\usepackage{setspace}
\usepackage{parskip}

\title{Before the Is: On the Normative Conditions of Knowing}
\author{}
\date{}

\begin{document}

\maketitle

\section*{Introduction}

A familiar structure in moral philosophy begins with a distinction between what \textit{is} and what \textit{ought} to be. We establish, as far as possible, what the case is, and then ask what follows normatively from that state of affairs.

I would like to suggest that this order is incomplete.

The problem is not only how we move from \textbf{IS} to \textbf{OUGHT}. The more prior problem is how a finite knower can establish a sufficiently reliable \textbf{IS} in the first place.

An \textit{is} available to a human being is never simply the world as it exists independently of the knower. It is the output of a model: a structured representation formed from available information, interpreted through a particular perspective, and necessarily limited by what that perspective can access.

This creates a prior normative question:

\begin{quote}
    \textbf{What must a model do in order to remain a model of reality rather than a model of itself?}
\end{quote}

My suggestion is that the answer introduces normativity before the conventional is/ought transition. A model that claims to represent reality must remain open to relevant information that could contradict and modify its present representation. If it systematically excludes information because that information threatens the coherence of the model, the model may preserve its internal consistency, but it progressively ceases to model what is outside itself. It begins to model the conditions of its own persistence.

The resulting claim is therefore stronger than the familiar observation that our knowledge is fallible:

\begin{quote}
    \textbf{A reliable IS requires a prior OUGHT: a norm governing the model's access to information and its openness to correction.}
\end{quote}

And because, for finite human beings, relevant information is often located in perspectives other than our own, this epistemic norm has a relational consequence.

The conditions for knowing what is include conditions for remaining in relation with sources of information that can show us what our own model cannot see.

This may provide an unexpected bridge between epistemology, ethics, and the philosophy of relationship.

\section*{The apparent priority of IS}

The classical form of the problem is familiar:

\[
\text{IS} \rightarrow \text{OUGHT}
\]

We first determine what is the case. We then determine what ought to be done.

The distinction is indispensable. We should not simply smuggle normative conclusions into descriptions of the world.

But the distinction can obscure a prior problem.

\textbf{What qualifies as the IS from which the OUGHT is to be derived?}

For an omniscient observer, the question might be trivial. For a finite observer, it is not.

A human being does not encounter the totality of reality. We encounter information. We select, organize, interpret, remember, compare, and integrate that information into models.

What we subsequently call \textit{the situation} is therefore already the result of a process.

This does not imply that reality is constructed by the observer. It means something more modest and more consequential:

\begin{quote}
    \textbf{Our access to what is always occurs through a model of what is.}
\end{quote}

The model may be more or less accurate. It may be corrected. It may be radically wrong. But it is the mechanism through which a finite knower can have an operational representation of reality at all.

The question of how the model should function therefore precedes the question of what follows from its output.

\section*{When does a model cease to model reality?}

Consider a model (\(M\)) receiving information (\(I\)).

When the information is consistent with the model, nothing fundamental is required. But suppose the incoming information contradicts the model:

\[
I \not\approx M
\]

There are two broad possibilities.

The model can change:

\[
M + I \rightarrow M'
\]

Or it can exclude the information:

\[
M + I \rightarrow M
\]

The second response can be psychologically and computationally attractive. A model that excludes contradictory information protects its existing coherence.

But something important has happened.

The model is no longer allowing information from outside itself to determine whether it should change.

It has become self-protective.

This does not prove that the model is false. A model can reject information because that information is genuinely false, irrelevant, or unreliable. The problem arises when information is rejected \textbf{because it conflicts with the model itself}.

At that point, the model contains a rule of the form:

\begin{quote}
    Information that contradicts this model shall not be allowed to change this model.
\end{quote}

Such a rule makes the model increasingly self-referential.

The model no longer primarily asks:

\begin{quote}
    \textit{What does the information tell me about reality?}
\end{quote}

It begins to ask:

\begin{quote}
    \textit{What must I do with the information so that my existing model can remain intact?}
\end{quote}

That is the transition I want to mark.

\begin{quote}
    \textbf{A model can remain internally coherent while losing its capacity to model what lies outside itself.}
\end{quote}

The danger of closure is therefore not simply error. It is the loss of the mechanism by which error can become visible.

\section*{The first OUGHT}

This produces a normative requirement before the familiar moral OUGHT.

If a system intends to maintain a reliable representation of reality, it must preserve the possibility that relevant information can alter that representation.

In schematic form:

\[
\text{OUGHT}_0 \rightarrow \text{MODEL} \rightarrow \text{IS}
\]

The first OUGHT is not yet a moral command such as \textit{help}, \textit{do not harm}, or \textit{remain married}.

It is a condition of epistemic functioning:

\begin{quote}
    \textbf{A model that purports to represent reality ought not to exclude relevant information merely because that information threatens the model's existing representation.}
\end{quote}

This is a normative statement because the model has alternatives. It can be open or closed. It can privilege correction or self-preservation.

The crucial point is that the norm cannot simply be derived from the resulting IS. It is a condition under which the IS can become sufficiently reliable to serve as the premise for subsequent normative reasoning.

Thus the familiar structure becomes:

\[
\text{OUGHT}_0 \rightarrow \text{MODEL} \rightarrow \text{IS} \rightarrow \text{OUGHT}_1
\]

where (\(\text{OUGHT}_0\)) governs the formation of a sufficiently reliable IS, and (\(\text{OUGHT}_1\)) concerns what ought to be done given that IS.

This does not abolish the is/ought distinction.

It introduces a \textbf{prior layer of normativity}.

\section*{There is no need to deny an objective reality}

This argument does not require the claim that there is no fixed reality.

There may be a perfectly determinate reality independent of every observer.

The weaker and more defensible claim is sufficient:

\begin{quote}
    \textbf{For a finite knower, there is no completely perspective-independent access to that reality.}
\end{quote}

The world may be fixed.

Our access to it is not.

Consequently, what appears as \textit{IS} to a finite system is always provisional with respect to the information available to that system.

This distinction matters.

The claim is not:

\begin{quote}
    There is no truth.
\end{quote}

It is:

\begin{quote}
    \textbf{No finite model can assume that its current representation exhausts the truth without thereby making an epistemic error about its own limitations.}
\end{quote}

A model that knows itself to be limited must therefore preserve the possibility of information that lies beyond its present state.

\section*{The relational consequence}

At this point, relationship enters the argument without requiring any initial appeal to morality.

A finite knower cannot generate all relevant information from within itself.

Other people occupy different positions. They have different histories, observations, memories, sensitivities, measurements, and experiences.

The other therefore represents something structurally important:

\begin{quote}
    \textbf{information that is not available from my perspective alone.}
\end{quote}

This makes relationship an epistemic structure.

At its minimum, a relationship is an information channel through which the state of one system can become available to another.

The relationship does not guarantee truth. The other person may be mistaken. But the existence of an independent perspective creates the possibility of information that my model could not have generated from itself.

This yields a form of reciprocal epistemic dependence:

\[
M_A \leftrightarrow M_B
\]

Neither model necessarily contains the truth. Each may contain information inaccessible to the other.

If both models remain open to information from the other, each can potentially reveal limitations in the other.

The relationship therefore becomes a condition for mutual correction.

And the circle closes:

\begin{quote}
    \textbf{My model needs the other precisely because my model is limited. The other needs me for precisely the same reason.}
\end{quote}

\section*{The paradox of an own reality}

This has an unexpected consequence for the idea of an ``own reality.''

I can legitimately experience the world from my own perspective. I can have memories, convictions, interpretations, and experiences that belong to me.

But if my reality is to remain a model of something beyond itself, it must contain the possibility that it is incomplete.

In other words:

\begin{quote}
    \textbf{My own reality must contain the possibility of being wrong about my own reality.}
\end{quote}

This is more than ordinary humility.

It is structurally necessary for an open model.

A completely closed ``own reality'' cannot be corrected by anything outside itself. But precisely for that reason, it cannot distinguish between \textit{being right} and \textit{being unable to discover that it is wrong}.

The paradox can therefore be stated as follows:

\begin{quote}
    \textbf{The freedom to have one's own reality requires the freedom to question it. But a reality that genuinely contains the question of its own validity is no longer a complete certainty.}
\end{quote}

The question mark is therefore not an external attack on the model.

It must be part of the model itself.

\section*{A concrete human case}

Consider two people who have shared many years of life and subsequently develop radically incompatible accounts of that shared history.

One person experiences the other as loving and safe. The other experiences the relationship as threatening and harmful.

From within their respective models, each may experience their own account as entirely real.

The temptation is to formulate the problem as a disagreement over conclusions:

\begin{quote}
    \textit{Given that this is what happened, what should we do?}
\end{quote}

But the deeper problem precedes that question.

The two people disagree about the IS itself.

One model contains:

\[
\text{IS}_A = \text{``safe and loving''}
\]

and therefore reaches one practical conclusion.

The other contains:

\[
\text{IS}_B = \text{``dangerous and harmful''}
\]

and therefore reaches another.

The important question is therefore not immediately which OUGHT follows.

It is:

\begin{quote}
    \textbf{What conditions would allow either model to establish whether its representation of the shared reality is incomplete?}
\end{quote}

This question cannot be answered by simply granting either person authority over their own model.

Nor does it follow that one person must be believed merely because their experience is sincere.

Sincerity is not the same as accuracy.

Instead, the problem requires access to information capable of challenging both models.

The disagreement therefore reveals, rather than resolves, the epistemic dependence between the two perspectives.

\section*{When the relationship itself becomes the casualty}

There is a further paradox.

Suppose two people end their relationship because their realities have become incompatible.

Ending the relationship may be entirely justified. There are circumstances in which maintaining a relationship is harmful or impossible.

But epistemically, something else happens at the same time.

The relationship may have been one of the few channels through which the two models could access information about their shared history that neither possessed independently.

When that channel closes, the possibility of mutual correction is reduced.

The very disagreement that causes the relationship to break may therefore also eliminate one of the mechanisms through which the disagreement could have been investigated.

The competing realities can become consolidated.

This does not mean that relationships must always be preserved.

It means that \textbf{the epistemic value of relationship is independent of whether the relationship should ultimately continue}.

The loss of a relationship can therefore be simultaneously:
\begin{itemize}
    \item the loss of a social bond,
    \item the loss of an information channel,
    \item and the loss of access to perspectives capable of correcting one's model.
\end{itemize}

This gives a deeper meaning to the proposition:

\begin{quote}
    \textbf{Maintaining a relationship can mean maintaining access to reality.}
\end{quote}

\section*{From epistemic norm to relational norm}

We can now return to the normative question.

If reliable representation requires openness to relevant information, and if relevant information is distributed across perspectives, then a norm concerning information access has relational implications.

The first norm is epistemic:

\begin{quote}
    \textbf{Do not close the model against relevant information merely because it threatens the model.}
\end{quote}

The relational implication is:

\begin{quote}
    \textbf{Do not unnecessarily close the channels through which relevant information can reach the model.}
\end{quote}

This still does not imply that every relationship must be maintained.

Nor does it imply that every person must be believed.

It implies something more precise:

\begin{quote}
    \textbf{A system that wishes to remain corrigible must preserve access to independent sources of information wherever those sources remain relevant to the question being modeled.}
\end{quote}

Relationship is therefore not intrinsically sacred because it is relationship.

Its epistemic significance derives from what it makes possible:

\textbf{access to information that the model cannot generate by itself.}

\section*{The deeper reversal}

We can now see why the ordinary formulation of the is/ought problem may be incomplete.

The familiar question is:

\begin{quote}
    \textbf{Given what is, what ought we to do?}
\end{quote}

But a prior question is required:

\begin{quote}
    \textbf{What ought we do so that what we take to be is sufficiently open to correction to serve as a reliable basis for action?}
\end{quote}

This is not a rejection of the is/ought distinction.

It is a distinction between two different normative levels.

The first concerns the \textbf{formation of knowledge}.

The second concerns \textbf{action on the basis of that knowledge}.

Thus:

\[
\text{OUGHT}_{\text{epistemic}} \rightarrow \text{MODEL} \rightarrow \text{IS} \rightarrow \text{OUGHT}_{\text{practical}}
\]

The first OUGHT is not an arbitrary moral preference. It is the condition under which the model can continue to function as a model of something beyond itself.

And because no finite model contains all relevant information, that condition is necessarily relational.

\section*{What this may mean for ethics}

This opens a possible line of inquiry.

Perhaps the foundations of ethics do not begin only with the question:

\begin{quote}
    \textit{What should I do, given the facts?}
\end{quote}

Perhaps they begin one level earlier:

\begin{quote}
    \textbf{What am I required to do to ensure that the facts on which I act remain open to correction?}
\end{quote}

If so, epistemic openness is not merely an intellectual virtue. It becomes a precondition for responsible action.

And this has a striking consequence.

The preservation of access to information may itself acquire normative significance.

Not because all information is true.

Not because every perspective deserves equal weight.

Not because every relationship must be preserved.

But because a finite agent that closes itself against relevant correction loses the capacity to distinguish between a model that is true and a model that merely protects itself from contradiction.

The fundamental ethical danger may therefore not be error alone.

It may be \textbf{epistemic closure}.

A closed model can be perfectly coherent. It can explain everything that it permits itself to see. It can even become extraordinarily certain.

But certainty is not the same as correspondence.

The deepest failure occurs when the model becomes so committed to preserving itself that information capable of changing it can no longer enter.

At that point:

\begin{quote}
    \textbf{The model ceases to function primarily as a model of the world and begins to function as a model of itself.}
\end{quote}

\section*{Conclusion}

The traditional is/ought distinction remains essential. We should not infer moral conclusions merely by describing the world.

But before asking what ought to follow from the \textit{is}, we may need to ask what makes an \textit{is} sufficiently reliable to support an ought at all.

For a finite knower, the answer appears to involve a prior norm:

\textbf{the model must remain open to relevant information capable of correcting it.}

Because relevant information is distributed across perspectives, this openness has a relational dimension.

The other is not merely someone who disagrees with my model.

The other is a possible source of information about what my model cannot see.

This creates a reciprocal structure:

\begin{quote}
    \textbf{I need the other because my model is limited. \\
    The other needs me because their model is limited. \\
    Our relationship creates access to information neither model possesses alone. \\
    That information makes correction possible. \\
    Correction makes a more reliable IS possible. \\
    And the IS then becomes the basis for further OUGHT.}
\end{quote}

The implication is therefore not that there is no objective reality.

It is that \textbf{our access to reality is conditional upon how we construct, maintain, and open our models}.

And this suggests a reversal worth taking seriously:

\begin{quote}
    \textbf{The problem is not only how we move from IS to OUGHT. The problem is that we first need an OUGHT governing the openness of our models before we can treat their output as a sufficiently reliable IS from which an OUGHT may responsibly follow.}
\end{quote}

If this is correct, then the ethical question begins earlier than we usually assume.

It begins with the question of whether we are keeping open the conditions under which reality can still correct us.

\end{document}